\documentclass{article}

% NeurIPS 2026 style (use [final] for camera-ready, [preprint] for arXiv)
\usepackage{neurips_2026}

\usepackage[utf8]{inputenc}
\usepackage[T1]{fontenc}
\usepackage{hyperref}
\usepackage{url}
\usepackage{booktabs}
\usepackage{multirow}
\usepackage{tabularx}
\usepackage{graphicx}
\usepackage{amsmath}
\usepackage{amssymb}
\usepackage{algorithm}
\usepackage{algpseudocode}
\usepackage{xspace}
\usepackage{xcolor}
\usepackage{caption}

\newcommand{\sys}{Wiki-Link Retrieval\xspace}
\emergencystretch=2em

\title{Wiki-Link Retrieval: Zero-Cost Inter-Document Linking\\for Curated Knowledge Bases}

\author{%
  Anonymous Author(s) \\
  Affiliation \\
  Address \\
  \texttt{email}
}

\begin{document}

\maketitle

\begin{abstract}
Most retrieval-augmented generation (RAG) systems treat documents as independent
units, ignoring the explicit link structure that already exists in curated
knowledge bases. GraphRAG-style approaches attempt to recover inter-document
relationships via LLM-based entity extraction, but this adds 10--100$\times$ indexing
cost and introduces extraction noise. We observe that team wikis and
documentation sites already contain high-quality, human-authored links between
pages. We propose \textbf{Wiki-Link Retrieval}: parsing existing wiki link structure
into an adjacency list and using it for post-retrieval 1-hop expansion after
standard ANNS search. On a controlled 20-document corpus, link expansion
improves Recall@5 by +3.5\% overall and +8.3\% on hard-link queries where the
relevant document has low lexical overlap. Wiki-link retrieval achieves 2.2$\times$
better Recall@5 and 3.2$\times$ better MRR than a GraphRAG baseline with only ${\sim}1$\,ms
additional latency --- no entity extraction, no graph database, no extra LLM calls.
\end{abstract}


\section{Introduction}
\label{sec:introduction}

Retrieval-augmented generation (RAG) has become the standard approach for
grounding large language models in domain-specific knowledge. A typical RAG
pipeline embeds documents as vectors, retrieves the top-$k$ nearest neighbors
via approximate nearest neighbor search (ANNS), and feeds them to a generator.

However, this approach treats each document as an \textbf{isolated unit}. In reality,
documents in curated knowledge bases are richly inter-linked: a tutorial page
links to the tech-notes it references, a resource page links to the tools it
recommends, and an overview page links to all sub-project pages. These links
encode semantic relationships that are \textbf{invisible to vector similarity}.

GraphRAG and property graph approaches attempt to recover inter-document
relationships by running LLM-based entity extraction over raw text. While
powerful, this introduces significant costs:

\begin{itemize}
  \item \textbf{Index cost}: 10--100$\times$ more expensive than vanilla ANNS (LLM calls per document)
  \item \textbf{Extraction noise}: LLM-extracted entities and relations are error-prone
  \item \textbf{System complexity}: requires a graph database, community detection, and graph summarization infrastructure
\end{itemize}

We observe that \textbf{curated wikis already contain the link graph for free}.
Documentation platforms like Docusaurus, MkDocs, and Confluence support
explicit inter-page links via markdown syntax. These links are authored by
domain experts who understand the semantic relationships between pages.

Instead of extracting a graph from text, we \textbf{parse the existing link
structure} and use it for lightweight post-retrieval expansion:

\begin{center}
\begin{tabular}{l}
\hline
Query $\rightarrow$ ANNS retrieves top-$k$ documents \\
\quad $\rightarrow$ 1-hop link expansion via pre-built adjacency list \\
\quad $\rightarrow$ Score: primary results + decay-weighted linked context \\
\quad $\rightarrow$ Return augmented top-$k$ \\
\hline
\end{tabular}
\end{center}

Our contributions:
\begin{enumerate}
  \item \textbf{Landscape analysis} showing that no mainstream RAG system exploits curated wiki links at retrieval time (\S\ref{sec:related})
  \item \textbf{System design} for zero-cost link graph construction and post-retrieval expansion with bidirectional adjacency (\S\ref{sec:method})
  \item \textbf{Auto-sync pipeline} that keeps the link graph current with wiki edits at negligible maintenance cost (\S\ref{sec:autosync})
  \item \textbf{Empirical evaluation} demonstrating measurable recall improvement on hard-link queries and 2.2$\times$ advantage over GraphRAG (\S\ref{sec:experiments})
\end{enumerate}


\section{Related Work: The Inter-Linking Gap in RAG Systems}
\label{sec:related}

\subsection{Landscape of Current Approaches}

Table~\ref{tab:landscape} surveys how mainstream LLM-powered knowledge systems handle
inter-document linking. The key finding: \textbf{no existing system exploits curated
wiki links at the retrieval level}.

\begin{table}[t]
\caption{How mainstream RAG/wiki solutions handle inter-document linking}
\label{tab:landscape}
\small
\begin{tabularx}{\textwidth}{@{}lllll@{}}
\toprule
\textbf{Approach} & \textbf{Examples} & \textbf{Inter-Linking Strategy} & \textbf{Index Cost} & \textbf{Query Cost} \\
\midrule
Wiki + Vector Search & Notion AI, Dify, FastGPT & None at retrieval & Low: $O(n)$ embed & ANNS only \\
GraphRAG & Microsoft GraphRAG, LightRAG & Full KG: entities + relations & \textbf{High: 10--100$\times$} & ANNS + graph \\
Property Graph & LlamaIndex, Zep, Neo4j & (S,P,O) triples in graph DB & Medium: NER + triple & +50--200ms \\
Hierarchical Memory & Mem0, MemGPT & Working$\to$episodic$\to$semantic & Medium: multi-level & Multi-level \\
\textbf{Wiki-Link (ours)} & --- & Parse wiki links $\to$ adjacency $\to$ 1-hop & \textbf{Low: $O(n)$+$O(|E|)$} & \textbf{+${\sim}$1ms} \\
\bottomrule
\end{tabularx}
\end{table}

\subsection{Analysis of the Gap}
\label{sec:gap}

\textbf{Wiki + Vector Search} systems (Notion AI, Confluence AI, Dify, FastGPT)
represent the status quo. They index wiki pages as independent documents and
retrieve by vector similarity. While their UIs often display ``related pages''
sidebars derived from wiki links, these links are \textbf{never used at the
retrieval level}. A user querying ``how to handle GPU OOM'' will only get
documents with high lexical/semantic overlap --- even if a closely related page
is one click away in the wiki.

\textbf{GraphRAG}~\citep{graphrag2024} addresses this gap by extracting entity graphs
from raw text using LLMs. It builds community summaries at index time and uses
them for global query answering. While powerful for exploratory queries, it
has three drawbacks for curated knowledge bases: (1) indexing cost is 10--100$\times$
higher due to LLM calls, (2) extracted entities are noisy compared to
human-authored links, and (3) the system requires a graph database and
community detection infrastructure.

\textbf{Property Graph} approaches (LlamaIndex, Zep, Neo4j) extract structured
(subject, predicate, object) triples and store them in a graph database.
Queries use Cypher or GQL for precise relationship traversal. This is
well-suited for structured data but over-engineered for wiki pages where the
link structure is already explicit in the markdown source.

\textbf{Hierarchical Memory} systems (Mem0~\citep{mem0}, MemGPT~\citep{memgpt}) organize knowledge into
working/episodic/semantic tiers. While they model temporal dynamics, they do
not address inter-document linking --- a page about KV Cache optimization and a
page about NPU memory management are stored independently even if the wiki
explicitly links them.

\subsection{Our Position}

\sys occupies a unique niche: it exploits \textbf{human-curated}
link structure (unlike GraphRAG's noisy extraction) at \textbf{zero additional index
cost} (unlike property graphs' NER pipeline) with \textbf{negligible query
overhead} (unlike graph traversal). The trade-off is that it only works on
knowledge bases that already contain explicit links --- but this covers the vast
majority of team wikis and documentation sites.


\section{Method}
\label{sec:method}

\subsection{Link Graph Construction}
\label{sec:construction}

At ingestion time, we parse wiki markdown files for inter-page links
(e.g., \texttt{[text](../other-page.md)}). Each link is resolved to a canonical
source name (e.g., \texttt{wiki:tech-notes/kv-cache-optimization}) and stored as
document metadata.

At index build time, we construct a \textbf{bidirectional adjacency list}:
if document $A$ links to document $B$, both $A \to B$ and $B \to A$ are added as valid
expansion edges. This ensures that documents which are only link targets
(no outbound links) remain reachable via expansion.

\begin{algorithm}[t]
\caption{Link Graph Construction}
\label{alg:construction}
\begin{algorithmic}[1]
\Function{RebuildLinkGraph}{$\mathcal{D}$}
  \For{each document $D \in \mathcal{D}$}
    \For{each linked source $S \in D.\text{metadata}[\text{linked\_source\_names}]$}
      \State Add edge $D \to S$ \Comment{forward}
      \State Add edge $S \to D$ \Comment{reverse}
    \EndFor
  \EndFor
\EndFunction
\end{algorithmic}
\end{algorithm}

Graph construction is $O(|E|)$ where $E$ is the total number of links ---
typically a few hundred for a team wiki. No LLM calls, no entity extraction,
no graph database.

\subsection{Post-Retrieval 1-Hop Expansion}
\label{sec:expansion}

At query time, after ANNS retrieval and reranking produce the initial top-$k$
results, we perform 1-hop link expansion:

\begin{algorithm}[t]
\caption{Post-Retrieval 1-Hop Link Expansion}
\label{alg:expansion}
\begin{algorithmic}[1]
\Function{ExpandHitsWithLinks}{$\mathcal{R}$, $G$, $m$, $\alpha$}
  \State $\mathcal{R}' \gets \mathcal{R}$
  \State $c \gets 0$
  \For{each hit $H \in \mathcal{R}$}
    \For{each neighbor $N \in G[H.\text{document\_id}]$}
      \If{$N \notin \mathcal{R}'$}
        \State Add $N$ to $\mathcal{R}'$ with score $= H.\text{score} \times \alpha$
        \State $c \gets c + 1$
      \EndIf
      \If{$c \geq m$}
        \State \Return $\mathcal{R}'$
      \EndIf
    \EndFor
  \EndFor
  \State \Return $\mathcal{R}'$
\EndFunction
\end{algorithmic}
\end{algorithm}

\textbf{Key parameters:}
\begin{itemize}
  \item $m$ (\texttt{max\_expansion}): maximum number of linked documents to inject (default: 3)
  \item $\alpha$ (\texttt{decay}): score fraction passed to linked documents (default: 0.5)
\end{itemize}

\subsection{Optimal Decay Factor}
\label{sec:decay}

Our ablation study (\S\ref{sec:ablation}) reveals a counter-intuitive finding: \textbf{lower decay
($\alpha = 0.3$) outperforms higher decay ($\alpha \in [0.5, 1.0]$)}. With $\alpha = 0.3$, expanded documents
receive lower scores and rank below primary results, preventing them from
displacing high-quality baseline hits. With $\alpha = 1.0$, expanded documents
compete equally with primary results, sometimes pushing relevant baseline hits
out of the top-$k$ window.

\textbf{Recommendation:} Use $\alpha = 0.3$ for conservative expansion that augments
without disrupting.

\subsection{Automatic Knowledge Synchronization}
\label{sec:autosync}

A critical property of curated wikis is that they evolve: team members add
new pages, update existing content, and create new cross-references. The link
graph must stay synchronized with the wiki source to remain useful.

We implement a \textbf{zero-maintenance auto-sync pipeline}:

\begin{center}
\begin{tabular}{c}
\texttt{wiki push (git)} $\rightarrow$ \texttt{systemd timer (every 30 min)} \\
$\rightarrow$ \texttt{git pull} $\rightarrow$ \texttt{ingest\_wiki.py (idempotent upsert)} \\
$\rightarrow$ \texttt{rebuild link graph} $\rightarrow$ \texttt{KB live}
\end{tabular}
\end{center}

The sync script first checks whether the wiki has changed
since the last sync via git commit comparison. If unchanged, ingestion is
skipped entirely --- making the overhead of no-op syncs negligible. When changes
are detected, all documents are re-ingested via upsert (create new, update
existing), and the link graph is fully rebuilt.

This contrasts sharply with GraphRAG, where wiki updates require re-running
the expensive entity extraction pipeline over all documents.


\section{Experiments}
\label{sec:experiments}

\subsection{Experimental Setup}
\label{sec:setup}

\textbf{Corpus:} 20 interconnected wiki documents with deliberate hub/spoke
structure. Hub documents (tutorials, overviews) have high keyword overlap with
queries and many outbound links. Spoke documents (deep tech-notes) have
specific technical content and lower keyword overlap.

\textbf{Ground truth:} 20 queries across three difficulty levels:
\begin{itemize}
  \item \textbf{Easy} (5): direct keyword match
  \item \textbf{Medium} (5): multi-document, partial expansion benefit
  \item \textbf{Hard-link} (10): relevant document has low keyword overlap, reachable only via links from hub documents
\end{itemize}

\textbf{Graph statistics:} 20 nodes, 88 bidirectional edges, average degree 4.4.

\textbf{Baselines:}
\begin{itemize}
  \item Vanilla ANNS (token-overlap search, no expansion)
  \item GraphRAG (entity co-occurrence graph with 1-hop entity expansion)
\end{itemize}

\subsection{Main Results}
\label{sec:main}

\begin{table}[t]
\caption{\sys vs baselines on 20-document corpus ($k=5$)}
\label{tab:main}
\centering
\begin{tabular}{@{}lccc@{}}
\toprule
\textbf{Method} & \textbf{Recall@5} & \textbf{MRR} & \textbf{Latency} \\
\midrule
Vanilla ANNS (baseline) & 0.717 & 0.887 & 1.5ms \\
\textbf{Wiki-Link ($\alpha = 0.5$)} & \textbf{0.742} & 0.887 & 1.4ms \\
\textbf{Wiki-Link ($\alpha = 0.3$)} & \textbf{0.767} & 0.887 & 1.4ms \\
GraphRAG (1-hop entity) & 0.342 & 0.277 & 0.1ms \\
\bottomrule
\end{tabular}
\end{table}

Table~\ref{tab:main} summarizes the main results. Key findings:
\begin{itemize}
  \item Link expansion improves Recall@5 by \textbf{+3.5\%} ($\alpha=0.5$) or \textbf{+7.0\%} ($\alpha=0.3$) over baseline
  \item Wiki-link achieves \textbf{2.2$\times$ better Recall@5} and \textbf{3.2$\times$ better MRR} than GraphRAG on the same corpus
  \item Latency overhead is \textbf{negligible} (${\sim}$1ms for dict lookups)
\end{itemize}

\subsection{Hard-Link Query Analysis}
\label{sec:hardlink}

On the 10 hard-link queries (where relevant docs have low keyword overlap):

\begin{table}[h]
\caption{Hard-link query performance ($n=10$)}
\label{tab:hardlink}
\centering
\begin{tabular}{@{}lcc@{}}
\toprule
\textbf{Metric} & \textbf{Baseline} & \textbf{Wiki-Link ($\alpha=0.5$)} \\
\midrule
Recall@5 & 0.600 & 0.650 \\
$\Delta$ & --- & \textbf{+0.050} \\
\bottomrule
\end{tabular}
\end{table}

The largest single-query improvement: a query about ``how to apply for lab GPU access'' went from Recall@5=0.50 to 1.00 (+0.50), because link
expansion found the NPU setup tutorial via the cluster guide page.

\subsection{Ablation Studies}
\label{sec:ablation}

\begin{table}[t]
\caption{Ablation on decay factor $\alpha$ ($k=5$, $m=3$, bidirectional)}
\label{tab:abl_decay}
\centering
\begin{tabular}{@{}cccc@{}}
\toprule
$\alpha$ & \textbf{Recall@5} & \textbf{MRR} & \textbf{Interpretation} \\
\midrule
0.3 & \textbf{0.767} & 0.887 & Best: expanded docs rank low \\
0.5 & 0.742 & 0.887 & Good: balanced \\
0.7 & 0.742 & 0.875 & Expanded docs compete \\
1.0 & 0.717 & 0.863 & Worst: equals baseline \\
\bottomrule
\end{tabular}
\end{table}

\begin{table}[t]
\caption{Ablation on max expansion $m$ ($k=5$, $\alpha=0.5$, bidirectional)}
\label{tab:abl_maxexp}
\centering
\begin{tabular}{@{}ccc@{}}
\toprule
$m$ & \textbf{Recall@5} & \textbf{Note} \\
\midrule
1 & 0.742 & Same as higher values \\
3 & 0.742 & Budget consumed by seen neighbors \\
5 & 0.742 & No additional benefit \\
10 & 0.742 & Diminishing returns \\
\bottomrule
\end{tabular}
\end{table}

\begin{table}[t]
\caption{Ablation on link direction ($k=5$, $m=3$, $\alpha=0.5$)}
\label{tab:abl_dir}
\centering
\begin{tabular}{@{}ccc@{}}
\toprule
\textbf{Direction} & \textbf{Recall@5} & \textbf{Note} \\
\midrule
Outbound-only & 0.742 & Forward edges from hubs \\
Bidirectional & 0.742 & Same on this corpus \\
Baseline (no exp) & 0.717 & Reference \\
\bottomrule
\end{tabular}
\end{table}

\subsection{GraphRAG Comparison}
\label{sec:graphrag}

The GraphRAG baseline uses entity co-occurrence from regex NER (50 entities,
328 edges). Despite having 3.7$\times$ more edges than the wiki-link graph, it
achieves significantly worse retrieval quality because entity co-occurrence
is a noisy signal compared to human-authored links.

\begin{table}[t]
\caption{System complexity comparison}
\label{tab:complexity}
\centering
\begin{tabular}{@{}lcc@{}}
\toprule
\textbf{Dimension} & \textbf{GraphRAG} & \textbf{Wiki-Link} \\
\midrule
Graph nodes & 50 entities & 20 documents \\
Graph edges & 328 & 88 \\
Entity extraction & LLM/NER required & None \\
Graph database & Required & Not needed \\
Index build cost & High (LLM calls) & Low (link parse) \\
Human curation & None & Wiki links (existing) \\
\bottomrule
\end{tabular}
\end{table}


\section{Discussion}
\label{sec:discussion}

\subsection{When Does Link Expansion Help Most?}

Link expansion is most valuable when:
\begin{enumerate}
  \item \textbf{The relevant document has low lexical overlap} with the query but is linked from a keyword-rich hub document
  \item \textbf{The wiki has hub/spoke structure} with overview pages linking to detailed technical pages
  \item \textbf{The query is ``hard''} --- requiring context that goes beyond direct keyword matching
\end{enumerate}

\subsection{Limitations}

\begin{enumerate}
  \item \textbf{Requires existing link structure}: only works on wikis/docs with explicit inter-page links. Unlinked or poorly linked wikis won't benefit.
  \item \textbf{Expansion budget problem}: in dense graphs, most neighbors of a hit are already in the baseline result set, wasting the expansion budget.
  \item \textbf{Small corpus evaluation}: our controlled experiment uses 20 documents. Scaling to larger, mixed KBs (hundreds of non-wiki documents) requires further investigation --- on a 635-document KB with only 18 wiki pages, expansion showed $\Delta=0$ because wiki pages were already found by baseline.
\end{enumerate}

\subsection{Future Work}

\begin{itemize}
  \item \textbf{Smart expansion}: skip neighbors already in baseline without counting against the expansion budget
  \item \textbf{2-hop expansion}: neighbors-of-neighbors with deeper decay ($\alpha^2 = 0.09$)
  \item \textbf{Hybrid KBs}: combining wiki-link expansion with broader document retrieval in mixed-corpora settings
  \item \textbf{LLM-based GraphRAG comparison}: replace regex NER with actual LLM extraction for a fairer comparison
\end{itemize}


\section{Conclusion}
\label{sec:conclusion}

We presented \sys, a lightweight approach that exploits
human-authored link structure in curated wikis for post-retrieval expansion.
With zero additional indexing cost (no LLM calls, no graph database) and
negligible query overhead (${\sim}$1ms), it improves Recall@5 by up to +7\% and
outperforms GraphRAG by 2.2$\times$ on the same corpus. The approach fills a gap
in the current RAG landscape where no mainstream system exploits wiki links
at the retrieval level.


\begin{ack}
We thank the anonymous reviewers for their constructive feedback.
\end{ack}


\bibliographystyle{plainnat}
\bibliography{references}


\end{document}
